\section{関連研究}
\label{sec:related_works}

本節では，頻出パターンマイニングの既存手法を
(a) グローバル頻度型，(b) 時間制約型，(c) 局所密集型，
(d) トランザクション構造型の 4 軸で整理し，
各軸で未解決の問題を明確にする．
表~\ref{tab:related_comparison} にまとめを示す．

\subsection{グローバル頻度型手法}
\label{sec:related_global}

Apriori~\cite{agrawal1994fast} は反単調性に基づく候補生成・枝刈りの原型であり，
FP-Growth~\cite{han2000mining} は FP-tree による候補生成回避，
ECLAT~\cite{zaki2000scalable} は縦型表現（tidset）と集合積による効率化を達成した．
LCM~\cite{uno2004lcm} は頻出・極大・閉パターンの高効率列挙を実現する．
これらはいずれもデータ全体のグローバル支持度を評価するため，
局所密集区間を直接抽出する能力を持たない．
また共起単位は「1 トランザクション = 1 アイテム集合」であり，
同一トランザクション内に複数バスケットが存在する状況は想定されていない．

\subsection{時間制約型手法}
\label{sec:related_temporal}

PFPM~\cite{tanbeer2009discovering} は周期頻出パターンを対象とし，
最大ギャップ・最小周期長等の制約で探索を限定する．
PPFPM~\cite{kiran2017discovering} は partial periodicity を導入して
PFPM の柔軟性を改善したが，評価軸は依然として周期的出現頻度である．
LPFIM~\cite{mahanta2005finding} は局所頻出パターンを区間単位で扱う初期研究であるが，
区間抽出は gap 系しきい値に依存しパラメータ感度が残る．
LPPM~\cite{fournier2020local} は local periodic pattern を扱うが，
maxPer/maxSoPer/minDur の 3 パラメータへの依存が強い．

これらの時間制約型手法に共通する限界は，
密集性をgap/periodicity 制約に\emph{置き換える}点にある．
その結果，制約パラメータへの感度が高く，
ウィンドウ内局所支持度による密集の直接評価とは異なるアプローチとなる．

\subsection{局所密集型手法: DSAA2025 Apriori-window}
\label{sec:related_dsaa2025}

上記の縮退方向に対し，DSAA2025 の Apriori-window~\cite{dsaa2025} は，
固定幅スライディングウィンドウ上で局所 support を直接評価し，
gap/periodicity に問題を置き換えずに
局所密集そのものを一次問題として扱った．
これは密集パターン抽出における問題設定の質的転換であり，
パラメータ感度の問題を構造的に回避した点に意義がある．

ただし Apriori-window は入力を「1 トランザクション = 1 アイテム集合」と仮定するため，
同一トランザクション内に複数の独立バスケットが存在する状況では，
バスケット横断の誤共起（偽共起）を抑制できない．

\subsection{トランザクション構造と共起定義}
\label{sec:related_structure}

トランザクションの内部構造に着目した研究には複数の方向がある．
Inter-transaction association rules~\cite{lu2000inter} は
異なるトランザクション間の関係を扱うが，
同一トランザクション内のバスケット分離とは逆方向の拡張である．
Multi-level association rules~\cite{han1999mining} は
アイテム階層（例: 「飲料 $\supset$ 牛乳」）を導入して
異なる抽象度での頻出パターンを抽出するが，
同一時刻の複数購買イベントの分離は対象外である．

実務的には，POS データの日次集約や EC ログの
セッション単位集約において，同一顧客の複数購買や
複数顧客の同一時間帯購買が 1 レコードに混入する状況は一般的である．
しかし，この混入が頻出パターンマイニングの出力品質に与える影響を
定量的に分析し，アルゴリズム側で対処する研究は
我々の知る限り存在しない．

\subsection{本研究の位置づけ}
\label{sec:related_positioning}

本研究は，Apriori-window の核である「局所密集の厳密抽出」を維持しながら，
前手法群でも Apriori-window でも扱えなかった
バスケット構造由来の偽共起問題を対象化する．
具体的には，共起判定をトランザクション内からバスケット内へ拡張し，
DSAA2025 枠組みの上位互換として実運用データへの適用可能性を高める．

\begin{table}[t]
\caption{関連手法の比較（$\checkmark$: 対応，$\times$: 非対応）}
\label{tab:related_comparison}
\centering
\scriptsize
\begin{tabular}{lccc}
\toprule
手法 & 局所密集 & パラメータ非依存 & バスケット分離 \\
\midrule
Apriori / FP-Growth / ECLAT & $\times$ & --- & $\times$ \\
PFPM / PPFPM & $\triangle$ & $\times$ & $\times$ \\
LPFIM / LPPM & $\triangle$ & $\times$ & $\times$ \\
Apriori-window~\cite{dsaa2025} & $\checkmark$ & $\checkmark$ & $\times$ \\
\textbf{提案手法 (BAAW)} & $\checkmark$ & $\checkmark$ & $\checkmark$ \\
\bottomrule
\end{tabular}
\end{table}
